\begin{proposition}[Parabolic families and exceptional divisors]\hypertarget{tag-030E}{}\label{push:src:B.4.10}
\StatementRecord{push:src:B.4.10}{proved}{030E}{\SeriesRef{push:src:B.4.8}; \citet[Section 2, pp.1351--1353]{IstratiEtAl2022Root}; \citet[Introduction, pp.1750--1751]{Teleman2010AB}; \citet[Example 9.10, p.1386]{IstratiEtAl2022Root}; \citet[Proposition 5.6, pp.1369--1370]{IstratiEtAl2022Root}; \citet[Section 1.2]{DlousskyTeleman2012AB}}

Fix an integer $n\ge1$ and holomorphic units $a(s),\kappa(s)$ on a
disc, with $0<|a(s)|<1$. After shrinking the disc there is a proper
holomorphic submersion $\mathcal K\to\Delta_s$, with an embedded smooth
proper elliptic family
\[
 \mathcal E=(\C^*\times\Delta_s)/
       ((v,s)\sim(a(s)v,s)).
\]
Its central surface is a parabolic Inoue surface $S_n$, with
$\pi_1(S_n)=\Z$, $b_2(S_n)=n$, and elliptic axis of square $-n$.
For $s\ne0$, the surface $\mathcal K_s$ is obtained from the diagonal
Hopf surface $H_{\kappa(s)s^n,a(s)}$ by $n$ successive blow-ups at one
axis point and its indicated infinitely near strict-axis points.
In a fixed annulus marking of $\mathcal E$, its normal line has factor
\[
 g_s(v)=\kappa(s)(a(s)v+s)^n.
\]
This line extends across $s=0$ and is isomorphic to the holomorphic
family with factor $\kappa(s)a(s)^n v^n$.
For constant $a,\kappa$ it is a constant family of degree $-n$ lines.

The relation with the normal line of the unmodified Hopf axis is
\[
 N_{\mathcal E_s/\mathcal K_s}
   =L_{\kappa(s)s^n}\otimes\OO_{\mathcal E_s}(-n[-s]),
       \qquad s\ne0.
\]
Here $n[-s]$ is the total exceptional correction, with the
multiplicity of each successive blow-up retained. It has no
Cartier extension on $\mathcal E$, even after a finite base change or
the addition of a vertical component. The marked divisor used to change the normal line therefore
fails the extension condition in Proposition~\ref{push:src:B.4.8}.

For $n=8$, the degree $-8$ axis normal can be chosen to be the inverse
of $\OO_E(2)$ on any chosen smooth elliptic complete intersection
$E=(2,2)\subset\PP^3$. This supplies the normal-line matching with
$\operatorname{Bl}_E\PP^3$ described below, including one positive
triple section. The matching is the displayed isomorphism of the two normal lines,
with the stated triple-curve factor.
\end{proposition}


\begin{proof}
\begin{enumerate}[label=\arabic*.,leftmargin=*]
\item Let
$B_R\subset\C^2$ be a ball, and let
$\pi_n:\widehat B_R\to B_R$ be the following modification.
Blow up the origin, then successively the intersection of the newest
exceptional curve with the strict transform of the axis $x=0$.
In the last chart,
\[
 \pi_n(z,w)=(zw^n,w).
\]
Define the parameter-dependent embedding in that chart by
\[
 \sigma_s(z,w)=
       \bigl(\kappa(s)z,\ a(s)w+(1-a(s))s\bigr).
\]
Choose constants $A<\rho<1$ and $K>0$ with
$|a(s)|\le A$ and $|\kappa(s)|\le K$ on a closed smaller disc.
Shrink it further so that
$|a(s)w+(1-a(s))s|\le\rho R$ on $\overline B_R$.
Then
\[
 |\pi_n\sigma_s(z,w)|^2
 \le K^2R^2(\rho R)^{2n}+\rho^2R^2<R^2
\]
when $R$ is small enough. The last-chart coordinates of these images
are bounded, so this is compact containment of
$\sigma_s(\overline B_R)$ in $\widehat B_R$, uniformly in $s$.
For each finite bound on $|\kappa|$, these inequalities give such a radius $R$.

\item Remove the embedded ball and identify the outer and inner collars by
$\sigma_s\pi_n$. The outer collar avoids the exceptional locus, and
$\pi_n$ is a biholomorphism there; thus the collar identification is
a biholomorphism. Carry out the same construction over the parameter
disc. The embeddings and the identifications are holomorphic and
preserve $s$. Uniformly disjoint collars give a Hausdorff smooth
threefold, and the projection is a holomorphic submersion in every
chart. Uniform compact fundamental regions over compact parameter
sets prove properness. This is the usual Kato construction, with
the explicit moving embedding above; see
\citet[Section 2, pp.1351--1353]{IstratiEtAl2022Root}.
The general occurrence of deformations from Kato surfaces to
blown-up Hopf surfaces is classical; see
\citet[Introduction, pp.1750--1751]{Teleman2010AB}.
We use explicit formulas here to retain the elliptic marking and
its exceptional correction.

\item The contracted map $\pi_n\sigma_s$ fixes $(0,s)$.
In the coordinate $v=w-s$ it becomes
\[
 f_s(x,v)=
    \bigl(\kappa(s)x(a(s)v+s)^n,\ a(s)v\bigr).
\]
At $s=0$ this is the parabolic monomial germ with exponent matrix
$\left(\begin{smallmatrix}1&n\\0&1\end{smallmatrix}\right)$.
The chosen modification is exactly the $n$-fold modification in
\citet[Example 9.10, p.1386]{IstratiEtAl2022Root}.
Consequently the central surface is the stated parabolic Inoue
surface; its elliptic axis is disjoint from the rational cycle.
The construction gives $\pi_1=\Z$ and $b_2=n$, or equivalently one
can use the cited description of this surface.

\item The strict axis is preserved in the collar gluing. Its coordinate
$v$ is multiplied by $a(s)$. To see its full quotient, the axis disc
in the blowdown coordinates is
\[
 D_s=\{v:|v+s|<R\}.
\]
For small $s$, it contains $0$ and satisfies $a(s)D_s\Subset D_s$.
Its dilates have union $\C$ and shrinking intersection $\{0\}$.
Thus $D_s\setminus a(s)D_s$, with the two boundary collars glued,
is a fundamental annulus for $\C^*/a(s)^\Z$.
This proves the exact elliptic marking throughout the parameter disc.

\item There are two normal frames, whose distinction is essential.
In the last blow-up chart the collar map $\sigma_s\pi_n$ has factor
\[
 g_z(v)=\kappa(s)(v+s)^n.
\]
On a fixed annulus avoiding $v=0$, shrink $s$ so that both
$v+s$ and $a(s)v+s$ are nonzero. The blowdown normal coordinate is
$x=z(v+s)^n$, and in that frame the factor is
\[
 g_x(v)=\kappa(s)(a(s)v+s)^n.
\]
The change between the frames is a nonzero gauge on these annuli.
The central factor is $\kappa(0)a(0)^n v^n$ and has degree $-n$,
with the factor-of-automorphy convention of
\citet[Proposition 5.6, pp.1369--1370]{IstratiEtAl2022Root}.

\item An explicit regular gauge verifies the whole family. On a fixed
annulus and its overlap with its $a(s)$-translate put
\[
 H_s(v)=\prod_{j\ge0}\left(1+\frac{a(s)^j s}{v}\right)^n.
\]
The product converges locally uniformly in $(v,s)$ and is nonzero
after shrinking the parameter disc. It satisfies
\[
 \frac{H_s(a(s)v)}{H_s(v)}
      =\left(1+\frac{s}{a(s)v}\right)^n.
\]
Replacing $x$ by $x/H_s(v)$ changes $g_x$ to
$\kappa(s)a(s)^n v^n$.
The collar maps are linear in the normal coordinate, so the same
calculation identifies a tubular line neighbourhood.

\item For completeness, the generic Hopf blowdown also has an explicit
developing map. Fix $s\ne0$ and set
\[
 G_s(v)=\prod_{j\ge1}\left(1+\frac{a(s)^j v}{s}\right)^n,
       \qquad
 \Psi_s(x,v)=(xG_s(v),v).
\]
This product is entire in $v$, converges locally uniformly, and obeys
\[
 G_s(a(s)v)=\frac{G_s(v)}{(1+a(s)v/s)^n},\qquad
 \Psi_s f_s=\operatorname{diag}(\kappa(s)s^n,a(s))\Psi_s.
\]
For sufficiently small nonzero $s$, both diagonal multipliers have
absolute value less than one. Compose $\Psi_s$ with $\pi_n$ on a
Kato fundamental region. The displayed equivariance glues it to a
holomorphic map to $H_{\kappa(s)s^n,a(s)}$.
It never reaches the excluded origin: that would require
$x=v=0$, which is the fixed point $(0,s)$ removed with the inner ball.
All zero-factor curves of $G_s$ have
$v=-s/a(s)^j$, so they lie in the one axis orbit $[-s]$.
The original modification also contracts over $v=-s$.
At that representative,
\[
 G_s(-s)=\prod_{j\ge1}(1-a(s)^j)^n\ne0.
\]
Hence locally the map is precisely $\pi_n$ followed by a coordinate
unit. Away from that orbit it is a local biholomorphism of degree
one. Compactness makes it a proper modification, and the displayed
local model identifies it as the stated $n$ successive blow-ups.
The developing map $\Psi_s$ is defined on the punctured parameter disc $s\ne0$.

\item Each blow-up of a point on the strict elliptic curve subtracts that
point from its normal line. Following all $n$ stages gives
\[
 N_s=L_{\kappa(s)s^n}\otimes\OO_{\mathcal E_s}(-n[-s]).
\]
The multiplicity $n$ records total exceptional transforms.
The reduced final exceptional locus meets the strict axis only once.
One may check the formula directly: the standard theta factor
$\theta_a(av)=-v^{-1}\theta_a(v)$ shows that $\OO([p])$
has factor $-p/v$. At $p=-s$ the tensor product above has factor
\[
 \kappa(s)s^n(v/s)^n=\kappa(s)v^n,
\]
which is equivalent to $\kappa(s)a(s)^n v^n$ by a nonzero monomial
gauge. This checks the signs and degrees independently of the
successive-blow-up formula.

\item The divisor $n[-s]$ cannot extend. For constant $a$, choose two
distinct classes $[v_1],[v_2]\in\C^*/a^\Z$.
The sequences $s_j=-a^jv_i$ tend to zero, whereas the corresponding
divisors are constantly $n[v_i]$ and have different limits in
$\operatorname{Sym}^n(E_a)$.
For varying $a(s)$, choose the holomorphic logarithm
$\tau_a(s)=\log a(s)/(2\pi i)$ on the disc.
The Abel--Jacobi coordinate of $n[-s]$ is
$n\log(-s)/(2\pi i)$ modulo $\Z+\tau_a(s)\Z$.
On a positive parameter loop it winds $n$ times around a primitive
cycle in the smooth elliptic family. An extending section over the
disc would have zero winding. A finite base change of order $N$
multiplies this nonzero integer by $N$.
After removing a vertical Cartier elliptic component, an extending
divisor would still have a proper horizontal part of
degree $n$, finite over the disc, defining a continuous
section of the relative symmetric power, which has just been ruled
out. Nonproper exceptional divisors in deformations of Kato surfaces
are a classical phenomenon, studied by
\citet[arXiv version~2, Section~1.2, PDF pp.~3--5]{DlousskyTeleman2012AB}; the calculation here
identifies the precise weighted divisor in this marked family.

\item This also explains the scope of ordinary elementary modifications.
If a degree-zero line $L_m$ from \ref{push:src:B.4.8} is changed by point blow-ups,
then its new normal is $L_m\otimes\OO(-D_{\rm exc})$.
If both that normal and the weighted Cartier correction extended on
a semistable good-reduction interface, their tensor product would
extend $L_m$, contradicting \ref{push:src:B.4.8}. In this Kato family the correction is exactly $D_{\rm exc}=n[-s]$,
whose nonextension was proved above.

\item Finally let $E=(2,2)\subset\PP^3$ be smooth and set
$L=\OO_E(2)$, of degree eight. Its normal in $\PP^3$ is
$L\oplus L$. For $A=\operatorname{Bl}_E\PP^3$, the exceptional
divisor is $D=E\times\PP^1$ and, in the lines convention,
\[
 N_{D/A}=L\boxtimes\OO_{\PP^1}(-1).
\]
A bare Hopf-product cap has normal $M\boxtimes\OO$ with
$\deg M=0$. The required triple-divisor line would then be
$(L\otimes M)^{-1}\boxtimes\OO(1)$, which has no nonzero section
because its elliptic degree is $-8$. Thus the required triple-divisor line has no effective section
on this interface.

\item For the parabolic surface with $n=8$, changing $\kappa$ varies the
flat factor through all of $\operatorname{Pic}^0(E)$:
constant factors are parametrized by $\C^*/a^\Z$.
Choose its axis normal to equal $L^{-1}$.
In $B=S_8\times\PP^1$ the divisor $D=E\times\PP^1$ then has normal
$L^{-1}\boxtimes\OO$. With triple section $T=E\times\{p\}$,
\[
 N_{D/A}\otimes N_{D/B}\otimes\OO_D(T)
   =(L\boxtimes\OO(-1))\otimes
     (L^{-1}\boxtimes\OO)\otimes(\OO\boxtimes\OO(1))
   \simeq\OO_D.
\]
This proves precisely the first-order matching claimed.
\end{enumerate}
\end{proof}
